% 作者题证的编辑转录副本，不是独立下载的上游原件。
% 来源：https://terrytao.wordpress.com/2011/12/16/254b-notes-3-quasirandom-groups-expansion-and-selbergs-316-theorem/

\paragraph{Definition and convolution convention.}
A finite group $G$ is $D$-quasirandom if every non-trivial irreducible unitary representation of $G$ has dimension at least $D$. Given two functions $f,g\in\ell^2(G)$, we can define the convolution $f*g\in\ell^2(G)$ by the formula
\[f*g(x):=\sum_{y\in G}f(y)g(y^{-1}x)=\sum_{y\in G}f(xy^{-1})g(y).\]
\begin{proposition}[3, Mixing inequality]
Let $G$ be a finite $D$-quasirandom group, and let $f,g\in\ell^2(G)$. If at least one of $f,g$ has mean zero, then
\[\|f*g\|_{\ell^2(G)}\le D^{-1/2}|G|^{1/2}\|f\|_{\ell^2(G)}\|g\|_{\ell^2(G)}.\]
\end{proposition}
\begin{proof}
By subtracting a constant from $f$ or $g$, we may assume that $f$ or $g$ both have mean zero. Observe that $f*g$ (being a superposition of right-translates of $f$) also has mean zero. Thus, we see that we may define an operator $T_g:\ell^2(G)_0\to\ell^2(G)_0$ by setting $T_gf:=f*g$. It thus suffices to show that the operator norm of $T_g$ is at most $D^{-1/2}|G|^{1/2}\|g\|_{\ell^2(G)}$.

Fix $g$. We can view $T_g$ as a $|G|-1\times|G|-1$ matrix. We apply the singular value decomposition to this matrix to obtain singular values
\[\sigma_1\ge\ldots\ge\sigma_{|G|-1}\ge0\]
of $T_g$, together with associated singular vectors. The operator norm of $T_g$ is the largest singular value $\sigma_1$. The operator $T_g^*T_g$ is then a self-adjoint operator (or matrix) with eigenvalues $\sigma_1^2,\ldots,\sigma_{|G|-1}^2$. In particular, we have
\[\operatorname{tr}T_g^*T_g=\sigma_1^2+\ldots+\sigma_{|G|-1}^2.\]
Now, a short computation shows that $T_g^*T_gf=f*g*\tilde g$, where $\tilde g(x):=\overline{g(x^{-1})}$, and (by embedding $\ell^2(G)_0$ in $\ell^2(G)$, and noting that $T_g^*T_g$ annihilates constants) the trace can computed as
\[\operatorname{tr}T_g^*T_g=|G|g*\tilde g(0)=|G|\|g\|_{\ell^2(G)}^2.\]
Thus, if $V$ is the eigenspace of $T_g^*T_g$ corresponding to the eigenvalue $\sigma_1^2$ (so that the dimension of $V$ is the multiplicity of $\sigma_1$), we have
\[\dim(V)\sigma_1^2\le|G|\|g\|_{\ell^2(G)}^2.\]
Now observe that if $\tau:G\to U(\ell^2(G)_0)$ is the left-regular representation (restricted to $\ell^2(G)_0$) then
\[T_g^*T_g\tau(h)f=\tau(h)T_g^*T_gf\]
for any $f\in\ell^2(G)_0$ and $h\in G$ (this is a special case of the associativity of convolution). In particular, we see that $V$ is invariant under $\tau$. Since $\tau$ has no non-trivial invariant vectors in $\ell^2(G)_0$, we conclude from quasirandomness that $V$ has dimension at least $D$, and the claim follows.
\end{proof}
